%% FullCourse_10Lecs -- Lecture 10, Topic 10.5: Case Study 4 --- AI as Research Collaborator
%% Source: ./archive_pre_split/Lec10_Case_Studies_v2.tex (section 5)
%% Pass B: layer in refinements from
%%         ../../MiniCourse_8hr/Slides/Module4_Agentic_AI/M4_T5_Case_DynamicMacro.tex
%% Pass B ports: the canonical five-stage workflow vocabulary
%%         (model -> equilibrium -> algorithm -> pseudo-code -> code)
%%         as a reference framework before the labor-market case stages.

%% Shared preamble for the FullCourse_10Lecs topic decks.
%%
%% Each topic .tex \input's this file, then defines its own \title, \date
%% override (if any), \graphicspath, and \begin{document}.
%%
%% Reconstructed verbatim from the inlined copy preserved in
%% Lec10_Case_Studies/Lec10_T8_Case6_DeepRamsey.tex (lines 23-190), which is
%% explicitly annotated there as "from ../shared_preamble.tex, copied verbatim".
%% Base preamble matches MiniCourse_8hr/Slides/shared_preamble.tex; this version
%% additionally provides \sectiondivider, the conceptbox/cautionbox/examplebox/
%% takeawaybox tcolorboxes, and \compacttables used across the split decks.

\documentclass[10pt,english,aspectratio=169]{beamer}

\usepackage{ae,aecompl}
\usepackage[T1]{fontenc}
\usepackage{textcomp}
\usepackage[utf8]{inputenc}
\usepackage{babel}
\usepackage{datetime2}

\setcounter{secnumdepth}{3}
\setcounter{tocdepth}{3}

\usepackage{amsmath, amssymb, amsthm, graphicx}
\usepackage{booktabs, multirow}
\usepackage{tikz}
\usepackage{pgfplots}
\pgfplotsset{compat=1.16}
\usepackage[most]{tcolorbox}
\tcbuselibrary{skins, breakable, theorems}
\usepackage{listings}
\usepackage{xcolor}

\lstset{
  basicstyle=\ttfamily\small,
  keywordstyle=\color{blue!80!black}\bfseries,
  commentstyle=\color{green!50!black}\itshape,
  stringstyle=\color{orange!80!black},
  backgroundcolor=\color{gray!8},
  frame=single,
  rulecolor=\color{gray!40},
  breaklines=true,
  showstringspaces=false,
  numbers=none,
}

\lstdefinestyle{pythonstyle}{
    language=Python,
    basicstyle=\ttfamily\footnotesize,
    keywordstyle=\color{blue!80!black}\bfseries,
    commentstyle=\color{green!50!black}\itshape,
    stringstyle=\color{orange!80!black},
    frame=single,
    breaklines=true,
    showstringspaces=false,
    numbers=left,
    numbersep=5pt,
    numberstyle=\tiny\color{gray},
    backgroundcolor=\color{gray!8},
    rulecolor=\color{gray!40}
}

\ifx\hypersetup\undefined
  \AtBeginDocument{%
    \hypersetup{bookmarks=false,breaklinks=true,pdfborder={0 0 0},colorlinks=true,
      linkcolor=DarkText,urlcolor=RedTitle}%
  }
\else
  \hypersetup{bookmarks=false,breaklinks=true,pdfborder={0 0 0},colorlinks=true,
    linkcolor=DarkText,urlcolor=RedTitle}%
\fi

\makeatletter
\newcommand\makebeamertitle{%
  \begin{frame}[plain]
    \vspace*{1.0cm}
    \centering
    {\usebeamerfont{title}\usebeamercolor[fg]{title}\inserttitle\par}
    \vspace{1.0cm}
    {\usebeamerfont{author}\usebeamercolor[fg]{author}\insertauthor\par}
    \vspace{0.2cm}
    {\usebeamerfont{date}\usebeamercolor[fg]{date}\insertdate\par}
  \end{frame}}%
\AtBeginDocument{%
  \let\origtableofcontents=\tableofcontents
  \def\tableofcontents{\@ifnextchar[{\origtableofcontents}{\gobbletableofcontents}}
  \def\gobbletableofcontents#1{\origtableofcontents}%
}

\usetheme{metropolis}
\usefonttheme{professionalfonts}

\definecolor{LightGray}{RGB}{230,230,230}
\definecolor{RedTitle}{RGB}{0,0,200}
\definecolor{VeryLightGray}{RGB}{245,245,245}
\definecolor{DarkText}{RGB}{0,0,0}

\setbeamercolor{background canvas}{bg=white}
\setbeamercolor{title}{fg=RedTitle}
\setbeamercolor{author}{fg=DarkText}
\setbeamercolor{date}{fg=DarkText}
\setbeamercolor{frametitle}{bg=VeryLightGray,fg=DarkText}
\setbeamercolor{normal text}{fg=DarkText}
\setbeamerfont{title}{size=\large}

\setbeamersize{text margin left=5mm, text margin right=5mm}
\addtobeamertemplate{frametitle}{}{\vspace{-0.5em}}
\newcommand{\reducevspace}{\vspace{-0.5cm}}

% Shared section divider and box styles for split topic decks.
\newcommand{\sectiondivider}[2]{%
\begin{frame}[plain,noframenumbering]
\begin{center}
\vfill
{\Large\textbf{#1}\par}
\vspace{0.35cm}
{\normalsize #2\par}
\vfill
\end{center}
\end{frame}
}

\newtcolorbox{conceptbox}[1][]{
  colback=blue!3,
  colframe=RedTitle!55,
  boxrule=0.35mm,
  arc=1mm,
  left=2mm,
  right=2mm,
  top=1mm,
  bottom=1mm,
  #1
}
\newtcolorbox{cautionbox}[1][]{
  colback=red!3,
  colframe=red!50!black,
  boxrule=0.35mm,
  arc=1mm,
  left=2mm,
  right=2mm,
  top=1mm,
  bottom=1mm,
  #1
}
\newtcolorbox{examplebox}[1][]{
  colback=green!3,
  colframe=green!45!black,
  boxrule=0.35mm,
  arc=1mm,
  left=2mm,
  right=2mm,
  top=1mm,
  bottom=1mm,
  #1
}
\newtcolorbox{takeawaybox}[1][]{
  colback=VeryLightGray,
  colframe=RedTitle!55,
  boxrule=0.35mm,
  arc=1mm,
  left=2mm,
  right=2mm,
  top=1mm,
  bottom=1mm,
  #1
}
\newcommand{\compacttables}{\renewcommand{\arraystretch}{1.15}}

\setbeamertemplate{navigation symbols}{}
\setbeamertemplate{footline}{%
  \hfill%
  \usebeamercolor[fg]{page number in head/foot}%
  \usebeamerfont{page number in head/foot}%
  \insertframenumber\,/\,\inserttotalframenumber\kern1em\vskip2pt%
}

\makeatother

\author{Zhigang Feng}
% "date with time": compile timestamp so each build is identifiable.
\date{\today\ \DTMcurrenttime}


% Suppress redundant auto-generated metropolis section page; \sectiondivider frames are the dividers we keep.
\metroset{sectionpage=none}

\graphicspath{{./pic/}{./figures/}{./Exercise10_ES_Fellows/plots/}{../../MiniCourse_8hr/Slides/Module4_Agentic_AI/pic/}{../}}

\usetikzlibrary{positioning,arrows.meta,shapes,calc,fit,backgrounds}

\lstdefinestyle{markdownstyle}{
    basicstyle=\ttfamily\footnotesize,
    frame=single,
    rulecolor=\color{gray!40},
    backgroundcolor=\color{gray!8},
    breaklines=true,
    showstringspaces=false,
    numbers=none,
}

\title[AI for Econ Research]{AI for Economic Research: Dynamic Models, Language, and Agents\\
Lecture 10.5: Case Study 4 --- AI as Research Collaborator}

\begin{document}

\makebeamertitle

%=============================================================================
\begin{frame}{Topic 10.5 Agenda}
\begin{enumerate}
  \item \textbf{Reference framework} --- the canonical five-stage research workflow
  \item \textbf{Case Study 4} --- AI as research collaborator across the lifecycle
\end{enumerate}
\end{frame}

%=============================================================================
\section{Reference: The Five-Stage Workflow}
%===========================================================

\sectiondivider{Reference framework}{Model $\to$ equilibrium $\to$ algorithm $\to$ pseudo-code $\to$ code}

% ---- 5-Stage 1: Diagram -- ported from MiniCourse M4_T5 ----
\begin{frame}{The Canonical Five-Stage Workflow}
\begin{center}
\begin{tikzpicture}[
    box/.style={rectangle, draw=black, fill=VeryLightGray, rounded corners,
                minimum width=2.1cm, minimum height=0.9cm,
                text width=1.95cm, align=center, font=\footnotesize},
    arrow/.style={-{Stealth[length=2.5mm]}, thick}
]
    \node[box] (s1) {1\\\textbf{Model}};
    \node[box, right=0.35cm of s1] (s2) {2\\\textbf{Equilibrium}};
    \node[box, right=0.35cm of s2] (s3) {3\\\textbf{Algorithm}};
    \node[box, right=0.35cm of s3] (s4) {4\\\textbf{Pseudo-code}};
    \node[box, right=0.35cm of s4, fill=orange!15] (s5) {5\\\textbf{Code}};
    \draw[arrow] (s1)--(s2);
    \draw[arrow] (s2)--(s3);
    \draw[arrow] (s3)--(s4);
    \draw[arrow] (s4)--(s5);
\end{tikzpicture}
\end{center}

\medskip
\begin{itemize}\itemsep1pt
  \item \textbf{Stages 1--4} are \emph{contracts}. They stop the agent from inventing the problem while coding.
  \item \textbf{Stage 5} is execution. It is allowed only after the contracts are clear.
  \item \textbf{Human approval} happens between stages. Skipping a stage is a process failure, even if the code runs.
\end{itemize}
\end{frame}

% ---- 5-Stage 2: Each stage's artifact and gate ----
\begin{frame}{Each Stage Produces a Reviewable Artifact}
\small
\begin{tabular}{p{2.0cm} p{4.4cm} p{5.4cm}}
\toprule
\textbf{Stage} & \textbf{Artifact} & \textbf{Exit gate (human-checked)} \\
\midrule
1. Model & \texttt{delta\_spec.md}: equations, parameters, what stays unchanged & every new symbol has a defining condition \\
2. Equilibrium & policy, prices, distributions, aggregates, and clearing conditions & every new market has a clearing equation \\
3. Algorithm & method choice + numerical knobs + diagnostics & where it can fail; what convergence criterion \\
4. Pseudo-code & \texttt{pseudocode.md}: loop nesting, tensor shapes, validation gate & a 2nd-year PhD could implement without follow-up questions \\
5. Code & source files, tests, session notes & validation gate passes; reduce-to-previous behaves \\
\bottomrule
\end{tabular}

\medskip
\textbf{Mantra.} Stages 1--4 are \emph{contracts}; Stage 5 is \emph{execution}. The researcher's job is to design and approve the contracts.
\end{frame}

% ---- 5-Stage 3: Stage-specific prompt patterns ----
\begin{frame}[fragile]{Prompt Patterns for Stages 1--4}
\begin{lstlisting}[style=markdownstyle, basicstyle=\tiny\ttfamily]
# Stage 1 (Model) prompt pattern
Read versions/v{n}/model_spec.md.
Draft versions/v{n+1}/delta_spec.md for exactly one model extension.
Define equations, parameters, outputs, and what stays unchanged.
Restate the model in plain language; wait for my approval.

# Stage 2 (Equilibrium) prompt pattern
State the recursive equilibrium. List the policy, prices, distributions,
aggregates, and all clearing/optimality conditions. Push back on any
inconsistency in my model spec.

# Stage 3 (Algorithm) prompt pattern
Compare candidate methods and trade-offs. I am choosing {method}
because {reason}. Explain where it may fail, what diagnostics to log,
and what convergence criterion to use. Do not write code.

# Stage 4 (Pseudo-code) prompt pattern
Write Python-style pseudo-code: loop nesting, tensor shapes, loss terms,
diagnostics, and validation gate. Then have domain-reviewer check
Euler/KKT/market-clearing equations.
\end{lstlisting}
\end{frame}

% ---- 5-Stage 4: Stage 5 implementation and validation ----
\begin{frame}[fragile]{Stage 5 (Code): Implementation and Validation}
\textbf{Artifact.} Source files, tests, notebook section, and \texttt{session\_notes.md}.

\medskip
\begin{lstlisting}[style=markdownstyle, basicstyle=\tiny\ttfamily]
Implement V{n+1} per the approved pseudocode. Include diagnostics,
tests/test_v{n+1}_smoke.py, and the validation gate in delta_spec.md.
Do NOT add features beyond V{n+1}. Plan first; show me the plan
before writing code. Then run the tests and iterate until the gate
passes.
\end{lstlisting}

\medskip
\textbf{Agent loop.} Plan $\to$ approve $\to$ implement $\to$ verify $\to$ review $\to$ fix $\to$ score.

\medskip
\textbf{Stop condition.} The validation gate passes, \emph{and} reduce-to-previous checks behave as expected (e.g., $\psi_K=0$ recovers the prior version).

\medskip
\begin{tcolorbox}[colback=blue!3, colframe=RedTitle!50]\small
\textbf{Key lesson.} A passing test suite is not enough; the gates must encode \emph{economics} (procyclical capital, lifecycle hump, market clearing).
\end{tcolorbox}
\end{frame}

% ---- 5-Stage 5: Map the vocabulary onto Case 4 ----
\begin{frame}{Mapping the Five Stages onto Case Study 4 (Labor Market)}
\small
\begin{tabular}{p{2.2cm} p{4.6cm} p{4.8cm}}
\toprule
\textbf{Canonical stage} & \textbf{Case 4 instantiation} & \textbf{Artifact} \\
\midrule
1. Model & glue-work theory; AI-Easy vs.\ AI-Hard glue; $2\!\times\!2$ matrix & verbal theory + O*NET item codes \\
2. Equilibrium & empirical specification; allowed predictions P1, P2 & specification prompt + claims memo \\
3. Algorithm & 11-script pipeline architecture; \texttt{config.py} & pipeline DAG + 51 unit tests \\
4. Pseudo-code & per-script validation rules + adversarial review & module table + review memo \\
5. Code & scripts run; data executed; results interpreted & figures + repositioned paper \\
\bottomrule
\end{tabular}

\medskip
\textbf{Same vocabulary, different domain.} Stages 1--4 remain contracts; Stage 5 is execution. The discipline travels from dynamic macro solvers to empirical labor-market research \emph{intact}.
\end{frame}

%=============================================================================
\section{Case Study 4: AI as Research Collaborator}
%===========================================================

%===========================================================

\sectiondivider{Case Study 4}{AI as full-lifecycle research collaborator in empirical macroeconomics}

% ---- CS4-01: The Research Question ----
\begin{frame}{Case 4 --- The Research Question}
\begin{tcolorbox}[colback=blue!3, colframe=RedTitle!60]
\textbf{Question:} Why has rapid AI adoption not produced uniform occupational collapse,
and what role does coordinative labor play in shaping which occupations survive?
\end{tcolorbox}
\vspace{0.3cm}
\begin{itemize}
  \item Paper: ``Glue Work, Task Bundles, and the AI Labor Market'' --- Feng (2026, in progress)
  \item Extends Garicano--Li--Wu (2026) by endogenizing bundle strength via glue labor supply
  \item \textbf{This case covers the full research lifecycle:} theory articulation, empirical pipeline design, adversarial review, and claims discipline --- not just one execution stage
\end{itemize}
\end{frame}

% ---- CS4-02: The Glue Work Hypothesis ----
\begin{frame}{The Glue Work Hypothesis}
\begin{columns}[T]
\begin{column}{0.48\textwidth}
\textbf{What is glue work?}
\begin{itemize}
  \item Coordinative tasks that integrate components without producing a visible standalone output
  \item O*NET examples: coordinating work and activities; scheduling; communicating across silos
  \item Underrewarded: workers capture only $\lambda < 1$ of the social marginal product
  \item When glue erodes, the bundle weakens --- often with a lag before it becomes visible
\end{itemize}
\end{column}
\begin{column}{0.48\textwidth}
\textbf{The 2$\times$2 AI impact matrix}
\vspace{0.2cm}

\small
\begin{tabular}{p{2.0cm}p{1.6cm}p{1.6cm}}
\toprule
 & \textbf{Low AI glue} & \textbf{High AI glue} \\
\midrule
\textbf{High AI prod.} & Restructure & Full disruption \\
\textbf{Low AI prod.} & Stable & Glue erosion \\
\bottomrule
\end{tabular}
\vspace{0.3cm}

\footnotesize{Core prediction: high-AI + high-glue occupations \emph{restructure} rather than collapse.}
\end{column}
\end{columns}
\vspace{0.3cm}
\textbf{The insight was human. The matrix formalization was joint.}
\end{frame}

% ---- CS4-03: Stage 1 --- Theory Building ----
\begin{frame}{Stage 1 --- Theory Building with AI}
\begin{itemize}
  \item Human developed the core insight: endogenize bundle strength through glue labor supply
  \item Claude was asked to:
    \begin{itemize}
      \item Articulate the 2$\times$2 prediction matrix from verbal intuition
      \item Define AI-Easy Glue vs.\ AI-Hard Glue using specific O*NET item codes
      \item Identify the market failure formally ($\lambda < 1$ creates underinvestment equilibrium)
    \end{itemize}
  \item \textbf{AI-Easy Glue} (LLM-susceptible): documenting/recording information, synthesizing knowledge, routine cross-functional reporting
  \item \textbf{AI-Hard Glue} (embodied/relational): building trust, mentoring, navigating ambiguous authority, reading implicit organizational signals
  \item Prompting mode: \textbf{specification-driven} --- ``List O*NET items 4C1--4C3 that qualify as AI-Hard Glue by the following criteria\ldots'' not ``tell me what glue work is''
\end{itemize}
\end{frame}

% ---- CS4-04: Stage 2 --- Empirical Specification ----
\begin{frame}[fragile]{Stage 2 --- Translating Theory into Empirical Specification}
\begin{lstlisting}[style=markdownstyle, basicstyle=\tiny\ttfamily]
## Empirical Specification Prompt

You are translating the following theoretical model into a
testable empirical specification.

THEORY CLAIMS:
- P1: High-glue, high-AI-exposure occupations restructure
  rather than collapse (employment share positive 2019-2024)
- P2: The reallocation margin operates at junior entry
  (age < 30), not incumbent exit

VARIABLE CONSTRUCTION RULES:
- glue_core: mean of O*NET items [4C1.a.2, 4A4.a.1, 4A4.b.1]
  normalized to [0,1] at SOC 2018 level
- ai_exposure: FRS AIOE score (Felten-Raj-Seamans 2021)
- outcome: CPS monthly employment share (seasonally adj.)

FORBIDDEN: do not use occupation-level wage as a glue proxy
without first checking correlation > 0.6 with glue_core.
\end{lstlisting}
\footnotesize
\begin{itemize}
  \item Exact item codes prevent hallucinated variable definitions
  \item A ``FORBIDDEN'' instruction enforces scope before results appear
\end{itemize}
\end{frame}

% ---- CS4-05: Claims Discipline ----
\begin{frame}[fragile]{Claims Discipline --- Setting the Rules Before Seeing Results}
\begin{lstlisting}[style=markdownstyle, basicstyle=\tiny\ttfamily]
## Claims Discipline Specification

BEFORE any data is analyzed, record the following:

ALLOWED CLAIMS (theory-derived, testable):
- P1: Occupations in top quartile of both glue_core and AIOE
  show positive employment share growth 2019--2024
- P2: Entry share (age < 30) declines faster in
  low-glue / high-AI occupations

FORBIDDEN CLAIMS (data-mining temptations to block):
- Any claim about wage effects not pre-specified here
- Any interaction term not listed in the model section
- Any subgroup finding discovered post-estimation

GATE: results memo must cite which allowed claim each
regression speaks to before interpretation proceeds.
\end{lstlisting}
\footnotesize\textbf{This is claims discipline.} It must be set before results appear, not reconstructed after.
It is the empirical analog of pre-registration.
\end{frame}

% ---- CS4-06: Pipeline Architecture ----
\begin{frame}{Stage 3 --- The 11-Script Pipeline Architecture}
\begin{center}
\begin{tikzpicture}[
  pscript/.style={rectangle, rounded corners=3pt, draw=RedTitle!70, fill=blue!5,
    minimum width=1.42cm, minimum height=0.72cm, font=\scriptsize\bfseries,
    text width=1.38cm, align=center},
  cfgnode/.style={rectangle, rounded corners=3pt, draw=gray!60, fill=gray!8,
    minimum width=3.0cm, minimum height=0.48cm, font=\scriptsize, align=center},
  parr/.style={-{Stealth[length=2mm]}, thick, gray!60},
  darr/.style={dashed, gray!40, -{Stealth[length=1.5mm]}},
]
% Top row
\node[pscript] (s00) {00\\Download};
\node[pscript, right=0.22cm of s00] (s01) {01\\Glue\\Proxies};
\node[pscript, right=0.22cm of s01] (s02) {02\\Merge};
\node[pscript, right=0.22cm of s02] (s03) {03\\CPS\\Panel};
\node[pscript, right=0.22cm of s03] (s05) {05\\Groups};
% Config node
\node[cfgnode, above=0.5cm of s02] (cfg) {config.py (single source of truth)};
% Bottom row
\node[pscript, below=0.55cm of s00] (s06) {06\\Figures};
\node[pscript, right=0.22cm of s06] (s07) {07\\Regress.};
\node[pscript, right=0.22cm of s07] (s09) {09\\O*NET\\Test};
\node[pscript, right=0.22cm of s09] (s10) {10\\ASEC\\Panel};
\node[pscript, right=0.22cm of s10] (s11) {11\\Diagnos.};
% Top arrows
\draw[parr] (s00)--(s01); \draw[parr] (s01)--(s02);
\draw[parr] (s02)--(s03); \draw[parr] (s03)--(s05);
% Handoff
\draw[parr, very thick] (s05.south) -- ++(0,-0.27cm) -| (s06.north);
% Bottom arrows
\draw[parr] (s06)--(s07); \draw[parr] (s07)--(s09);
\draw[parr] (s09)--(s10); \draw[parr] (s10)--(s11);
% Config dashes
\foreach \t in {s00,s01,s02,s03,s05}
  \draw[darr] (cfg.south) -- (\t.north);
\end{tikzpicture}
\end{center}
\vspace{0.15cm}
\textbf{One config.py = single source of truth.} \quad 51 unit tests specified before any script runs on real data.
\end{frame}

% ---- CS4-07: Pipeline Module Detail ----
\begin{frame}{Pipeline Module Detail}
\small
\begin{tabular}{p{1.1cm}p{5.6cm}p{5.5cm}}
\toprule
Script & Key output & Validation rule \\
\midrule
00 & onet\_raw, frs\_aioe, crosswalks (parquet) & row count $\geq$ expected; no null SOC codes \\
01 & glue\_proxies (6 variants, SOC 2018) & all variants in $[0,1]$; correlation matrix saved \\
02 & master\_soc: glue + AI exposure + crosswalks & merge rate $\geq 95\%$; no duplicate SOC \\
03 & cps\_panel: employment share, entry share & weights sum to civilian pop; seasonal adj.\ flag set \\
05 & analysis: 4 AI $\times$ glue groups defined & group counts balanced; no SOC in $>1$ group \\
06--11 & figures, tables, diagnostics & files written; no NaN in reported statistics \\
\bottomrule
\end{tabular}
\vspace{0.3cm}
Each script has explicit inputs (parquet), explicit outputs (parquet), and a pass/fail validation rule. Failure halts the pipeline.
\end{frame}

% ---- CS4-08: Variable Construction ----
\begin{frame}{Stage 3 --- Variable Construction Details}
\begin{columns}[T]
\begin{column}{0.46\textwidth}
\textbf{Glue proxies (6 variants)}
\begin{itemize}\small
  \item \texttt{glue\_core}: 3 O*NET items, baseline
  \item \texttt{glue\_ext}: 8-item extended set
  \item \texttt{glue\_aihard}: AI-Hard items only
  \item \texttt{glue\_social}: social/interpersonal subset
  \item \texttt{glue\_dominance}: mean glue / mean all activities
  \item \texttt{glue\_minus\_mgmt}: robustness check
\end{itemize}
\vspace{0.2cm}
\footnotesize{The 6 variants test whether the theory's protection story depends on which aspect of coordination is measured.}
\end{column}
\begin{column}{0.50\textwidth}
\textbf{AI exposure and outcomes}
\begin{itemize}\small
  \item AI exposure: FRS AIOE (2021), SOC-linked
  \item Outcomes: CPS monthly employment share
  \item Entry margin: age $<30$ share within occ-month
  \item Long-run: ASEC panel 2000--2025
  \item 4 groups: median splits of AIOE $\times$ glue\_core
\end{itemize}
\vspace{0.2cm}
\textbf{Core group definitions:}

\small
\begin{tabular}{ll}
High AI + High Glue & Restructuring \\
High AI + Low Glue  & Disruption \\
Low AI + High Glue  & Insulated \\
Low AI + Low Glue   & Stable \\
\end{tabular}
\end{column}
\end{columns}
\end{frame}

% ---- CS4-09: Adversarial Review ----
\begin{frame}{Stage 4 --- Adversarial Review}
\begin{tcolorbox}[colback=blue!3, colframe=RedTitle!60]
\textbf{Adversarial prompt:} ``Find 6 structural weaknesses in this empirical design. For each, state whether it is fatal or manageable, and why.''
\end{tcolorbox}
\vspace{0.2cm}
Six weaknesses identified:
\begin{enumerate}\small
  \item AIOE is cross-sectional --- does not capture within-occupation AI adoption dynamics
  \item Glue proxies are task-level O*NET ratings, not worker-level observations
  \item Selection: high-glue workers may self-select into high-AI firms for unobserved reasons
  \item Crosswalk from SOC 2018 to Census codes loses $\sim$12\% of occupations (a validation rule catches this)
  \item CPS monthly sample sizes for narrow occupations produce large standard errors
  \item \textbf{The interaction (AI $\times$ Glue) is negative in the data} --- but the theory predicts positive
\end{enumerate}
\vspace{0.2cm}
\textbf{Item 6 was the decisive finding.} It repositioned the entire paper.
\end{frame}

% ---- CS4-10: When Data Contradicts Theory ----
\begin{frame}{When Data Contradicts Theory}
\begin{tcolorbox}[colback=blue!3, colframe=RedTitle!60]
\textbf{Data finding:} the estimated interaction (AI exposure $\times$ glue score) is \emph{negative}.\\
\textbf{Theory predicted:} positive --- glue preserves bundles, so high-glue/high-AI should survive.\\
\textbf{This is not a bug.} It is information.
\end{tcolorbox}
\vspace{0.2cm}
\begin{columns}[T]
\begin{column}{0.46\textwidth}
\textbf{What the theory predicted}
\begin{itemize}\small
  \item High-AI + High-Glue $\to$ restructuring and survival
  \item Glue preserves value when production tasks are automated
  \item Expected positive interaction
\end{itemize}
\end{column}
\begin{column}{0.50\textwidth}
\textbf{The reinterpretation}
\begin{itemize}\small
  \item Negative sign: high-glue workers in high-AI occupations are displaced, not preserved
  \item Effect operates through \emph{occupational reallocation}, not within-occupation survival
  \item Theory-specific prediction $\neq$ generic heterogeneity claim
\end{itemize}
\end{column}
\end{columns}
\vspace{0.2cm}
\textbf{Claude distinguished theory-specific predictions from generic heterogeneity findings.} That distinction determines what the paper can claim.
\end{frame}

% ---- CS4-11: Division of Labor ----
\begin{frame}{Division of Labor Across the Full Lifecycle}
\small
\begin{tabular}{p{3.4cm}p{3.2cm}p{5.0cm}}
\toprule
Task & Human Researcher & AI (Claude) \\
\midrule
Core insight & Endogenizes bundle strength via glue & \\
Literature positioning & Sets the GLW extension & \\
Variable specs & Co-designed, approved & Elaborated O*NET codes; caught crosswalk inconsistency \\
Pipeline architecture & Approved the design & Designed 11-script layout from scratch \\
Claims discipline & Enforces the gate & Drafted allowed/forbidden claims \\
Adversarial review & Decided which weaknesses were fatal & Structured the 6 adversarial findings \\
Data execution & Runs the scripts & \\
Interpretation & Final call on repositioning & \\
Paper writing & Authors & Assists with exposition \\
\bottomrule
\end{tabular}
\end{frame}

% ---- CS4-12: Key Prompting Strategies ----
\begin{frame}{Key Prompting Strategies in Case 4}
\begin{enumerate}
  \item \textbf{Specification-driven:} exact O*NET item codes, not vague requests --- prevents hallucinated variable definitions
  \item \textbf{Constraint-based (claims discipline):} allowed and forbidden claims specified before results appear --- prevents data mining
  \item \textbf{Adversarial:} ``find 6 structural weaknesses'' not ``review this memo'' --- prevents sycophantic validation
  \item \textbf{Modular decomposition:} each script as an independent testable unit --- prevents silent errors accumulating across stages
  \item \textbf{Deferred execution:} pipeline fully specified before any real data is analyzed --- separates design from execution, preserving pre-specification discipline
\end{enumerate}
\vspace{0.3cm}
\textbf{Each strategy is a structural response to a specific failure mode, not a stylistic preference.}
\end{frame}

% ---- CS4-13: Failure Modes ----
\begin{frame}{Failure Modes Specific to Case 4}
\begin{enumerate}
  \item \textbf{AI agrees too readily:} asking ``does this design look good?'' produces validation; adversarial mode must be explicit
  \item \textbf{Vague theory $\to$ untestable claims:} ``glue work matters'' is not a prediction; specification-driven prompting forces operationalization
  \item \textbf{Data mining disguised as exploration:} claims discipline must be logged before results; retroactive framing is indistinguishable from pre-specification to a reader
  \item \textbf{Silent crosswalk errors:} the 12\% SOC loss is found only because a validation rule checks merge rate; without validation rules, the pipeline produces credible-looking numbers
  \item \textbf{Mechanical overclaiming:} when data contradicts theory, the temptation is ad hoc reframing; adversarial review forces the repositioning to be deliberate and documented
\end{enumerate}
\end{frame}

% ---- CS4-14: The Meta-Irony ----
\begin{frame}{The Meta-Irony: AI Doing Glue Work}
\begin{columns}[T]
\begin{column}{0.46\textwidth}
\textbf{What the paper studies}
\begin{itemize}
  \item Glue work: coordinative labor that binds complex task bundles
  \item Holds specifications together across stages
  \item Tracks implicit commitments across time
  \item Underrewarded because its contribution is diffuse in output attribution
\end{itemize}
\end{column}
\begin{column}{0.50\textwidth}
\textbf{What Claude did in the research itself}
\begin{itemize}
  \item Maintained specification coherence across 11 scripts and 4 research stages
  \item Tracked which claims were allowed before results arrived
  \item Structured the adversarial review to produce actionable findings
  \item Held the interface between theory and empirical design
\end{itemize}
\end{column}
\end{columns}
\vspace{0.4cm}
\begin{center}
\textbf{The AI is demonstrating its own research subject.}\\
That is not an accident --- it is evidence about where AI comparative advantage lies in complex intellectual work.
\end{center}
\end{frame}

% ---- CS4-15: Key Lessons ----
\begin{frame}{Key Lessons from Case Study 4}
\begin{itemize}
  \item AI can operate across the \textbf{full research lifecycle} --- theory articulation, pipeline design, adversarial review --- not just execution
  \item The human researcher must own the core insight, literature positioning, data execution, and final interpretation
  \item \textbf{Specification-driven and adversarial prompting} are structural defenses against the most common AI failure modes in research settings
  \item \textbf{Claims discipline} must be established before results appear; it cannot be reconstructed credibly after
  \item When data contradicts theory, the right move is deliberate repositioning, not ad hoc reframing; AI can structure that repositioning but cannot decide whether it is intellectually honest
  \item This is a new pattern: \textbf{AI as full-lifecycle research collaborator}, distinct from skill, empirical pipeline, or dataset harness
\end{itemize}
\end{frame}


\end{document}
